%%%%%%%%%%%%%%%%%%%%%%%%%%%%%%%%%%%%%%%%%%%%%%%%%%%%%%%%%%%%%%%%%%%%%%%%%%%%%%
%% GLADE: Gap-aware Lightweight Adaptive Discretisation Engine for
%% Fuzzy-Pattern Tsetlin Machines on IoT/IoMT Intrusion Detection
%%
%% Target venue: IEEE Transactions on Industrial Informatics (TII)
%% Template    : IEEEtran v1.8 (two-column, 10pt)
%%%%%%%%%%%%%%%%%%%%%%%%%%%%%%%%%%%%%%%%%%%%%%%%%%%%%%%%%%%%%%%%%%%%%%%%%%%%%%
\documentclass[journal,10pt]{IEEEtran}

\usepackage{cite}
\usepackage{amsmath,amssymb,amsfonts}
\usepackage{algorithmic}
\usepackage{graphicx}
\usepackage{textcomp}
\usepackage{xcolor}
\usepackage{booktabs}
\usepackage{array}
\usepackage{url}
\usepackage{hyperref}
\usepackage{tikz}
\usetikzlibrary{arrows.meta,positioning,shapes.geometric,calc,fit}
\usepackage{multirow}
\usepackage{booktabs}
% \usepackage{multirow}
\usepackage{graphicx}
\usepackage{subcaption}
\usepackage{longtable}

\def\BibTeX{{\rm B\kern-.05em{\sc i\kern-.025em b}\kern-.08em
    T\kern-.1667em\lower.7ex\hbox{E}\kern-.125emX}}

\begin{document}

\title{GLADE: A Gap-Aware Unsupervised Adaptive Binariser \\
       for Edge-Deployed Tsetlin Machines on \\
       IoT and IoMT Intrusion Detection}

\author{Rakesh~Reddy~Yakkati,~\IEEEmembership{Student~Member,~IEEE},
        and~Co-Authors%
\thanks{Manuscript submitted April 2026.}%
\thanks{R. R. Yakkati is with the Department of \textit{[department]},
\textit{[institution]} (e-mail: yrakesh2109@gmail.com).}}

\markboth{IEEE Transactions on Industrial Informatics,
          Vol.~XX, No.~X, Month~2026}%
         {Yakkati \MakeLowercase{\textit{et al.}}: GLADE for Edge-Deployed Tsetlin Machines}

\maketitle

%% ========================================================================
\begin{abstract}

\end{abstract}

\begin{IEEEkeywords}
Tsetlin machine, fuzzy pattern, unsupervised binarisation,
interpretable classification, intrusion detection, Internet of Things,
Internet of Medical Things, edge computing, TinyML.
\end{IEEEkeywords}



%% ========================================================================
\section{Introduction}
\IEEEPARstart{T}{he} proliferation of connected devices in
Internet-of-Things (IoT) and Internet-of-Medical-Things (IoMT)
deployments has created an operational need for intrusion detection to
be performed at the endpoint rather than in a remote gateway or cloud
service. On-device operation is motivated by the latency budgets of
bedside monitoring, by the privacy and regulatory constraints that
govern clinical telemetry, by the cost and unreliability of continuous
uplinks, and by the resilience benefits of keeping the decision loop
local when network connectivity is degraded. As a consequence, the
design space for such a detector is bounded not only by classification
accuracy but also by flash storage, random-access memory, per-sample
inference latency, and average power draw, so that model compactness
and on-device interpretability are raised to first-class design
constraints.

Tree-ensemble methods such as XGBoost~\cite{chen2016xgboost} and random
forests~\cite{breiman2001random}, together with feed-forward neural
networks trained through scikit-learn~\cite{pedregosa2011scikit}, are
routinely reported to dominate the public IoT intrusion-detection
leaderboards. However, their trained artefacts are typically measured
in megabytes, and their per-sample inference cost is not competitive
with an integer rule set on constrained hardware. The Tsetlin Machine
(TM) family, first introduced by Granmo~\cite{granmo2018tsetlin} and
subsequently extended with convolutional~\cite{granmo2019convolutional},
regression, coalesced and clause-sharing variants, has been studied as
an interpretable and hardware-efficient alternative in which every
class decision is expressed as a voting ensemble of conjunctive clauses
over Boolean literals. This representation yields a compact integer
artefact and a natural path to microcontroller deployment, as
demonstrated by REDRESS~\cite{maheshwari2023redress} on an STM32 device.
Further refinements, including sparse literal
budgets~\cite{bhattarai2023concise} and drop-clause
regularisation~\cite{sharma2023drop}, have progressively reduced the
TM memory footprint. The most recent member of the family, the
Fuzzy-Pattern Tsetlin Machine (FPTM) of Hnilov~\cite{hnilov2025fuzzy},
replaces the all-or-nothing clause rule of earlier TMs with a graded
fuzzy rule, and has been reported to reach comparable or superior
accuracy with between $50\times$ and $400\times$ fewer clauses.

A component of the TM pipeline that has attracted markedly less
attention than the machine itself is the \emph{binariser}: the
preprocessing stage in which continuous and categorical features are
converted into the Boolean literals on which every TM clause operates.
Because the information content of the literals lower-bounds the
attainable training quality, any inefficiency introduced at the
binarisation stage propagates directly into the trained clauses and
therefore into the final model footprint. Existing schemes such as the
\texttt{StandardBinarizer} shipped with the TMU library, the
equal-frequency \texttt{KBinsDiscretizer} of scikit-learn, and the
thermometer encoder, either assume a fixed bit width per feature,
depend on class labels, or disregard the extreme sparsity that is
characteristic of IoT traffic data. In the TON\_IoT
dataset~\cite{moustafa2021ton}, for example, more than two thirds of
the one-hot encoded columns contain above $80\%$ zero values after
preprocessing, which causes equal-frequency binning to collapse most
of the quantile edges onto the same value and to waste the entire bit
budget of that column.

\subsection{Contributions}
The central observation of this paper is that the choice of
binariser and the dynamics of fuzzy-pattern clause learning are not
independent: a binariser that allocates its bit budget adaptively
and that suppresses low-entropy thresholds shapes the literal
distribution in a way that is specifically favourable to a
downstream graded-vote learner. In other words, the learning
behaviour that emerges from the pairing is not the concatenation of
two unrelated stages but a joint design in which (i) the binariser
exposes informative literals with bounded marginal entropy, (ii)
the fuzzy clause evaluator absorbs the residual noise of those
literals through a graded rather than all-or-nothing decision, and
(iii) the two stages jointly produce short conjunctive rules whose
inference cost is dominated by a small number of 64-bit
bit-population counts. This joint design is the conceptual
contribution of the present paper; the individual stages are
separately well studied, but their behaviour \emph{together} is
new. The specific technical contributions are:

\begin{enumerate}
  \item An unsupervised adaptive binariser is proposed in which a
        per-feature bit budget is derived from the unique-value count
        and the zero fraction of each column, thresholds are placed by
        a conditional gap-aware procedure, low-entropy thresholds are
        discarded by a dead-bit elimination step, and the remaining
        thresholds are perturbed into high-variance neighbourhoods.
        Only a single integer hyperparameter is exposed to the user,
        and no cross-validation is required at fit time.
  \item A self-describing persistent storage format, the Template
        Binary Format (TBF), is introduced for the joint on-disk
        representation of the fitted binariser and the trained FPTM
        clauses. TBF is measured to reduce the on-disk model size by
        $40$--$60\%$ relative to JSON without altering inference
        semantics, and is designed to be readable by a portable
        $\sim$$150$-line C deserialiser.
  \item A strict single-core, fixed-hyperparameter benchmark is
        reported on four public IoT and IoMT datasets
        (NSL-KDD~\cite{tavallaee2009nslkdd},
        TON\_IoT~\cite{moustafa2021ton},
        WUSTL-EHMS-2020~\cite{hady2020wustl}, and MedSec-25) against
        twelve classical and neural baselines. The benchmark is
        accompanied by a binariser ablation, a GLADE-component
        ablation, a single-hyperparameter sensitivity study, a
        five-seed stability analysis, and a Wilcoxon signed-rank
        statistical comparison.
  \item An on-device deployment study on a Raspberry~Pi~5 is reported.
        Every trained model is shown to fit in at most $56$~kilobytes
        of on-disk storage. The GLADE+FPTM Numba JIT kernel predicts a
        sample in $4.10$--$7.42$~$\mu$s and a Numba-compiled DecisionTree
        reaches $1.37$--$1.39$~$\mu$s on a single Cortex-A76 core,
        matching the server-side accuracy exactly.
  \item The trained clauses are decoded into human-readable
        conjunctive rules over discretised feature literals, and are
        shown in several cases to reproduce the textbook description of
        the attack class they were trained to detect, without any
        post-hoc explainer.
\end{enumerate}

It is emphasised that the Fuzzy-Pattern Tsetlin Machine itself is
\textbf{not} claimed as a contribution of this work: it is used without
modification through the publicly released reference implementation of
Hnilov~\cite{hnilov2025fuzzy}. The novelty of the present paper is
concentrated on the binariser, on the persistent-storage format, and
on the demonstration that, once these two components are placed in
front of an unmodified FPTM, tree-ensemble accuracy can be matched on
IoT and IoMT intrusion-detection benchmarks at an on-disk budget that
is an order of magnitude smaller.

The remainder of the paper is organised as follows.
Section~\ref{sec:related} reviews the Tsetlin Machine family and the
intrusion-detection benchmarks used in the study.
Section~\ref{sec:binarisers} surveys the three binarisers in wide use
for TMs and presents the proposed GLADE binariser in detail.
Section~\ref{sec:fptm} describes the Fuzzy-Pattern Tsetlin Machine and
its bit-parallel evaluation. Section~\ref{sec:setup} defines the
experimental protocol, the preprocessing pipeline, and the twelve
baselines. Section~\ref{sec:results} reports the main benchmark and
all ablations. Section~\ref{sec:discussion} discusses limitations and
failure modes, and Section~\ref{sec:conclusion} concludes.

%% ========================================================================
\section{Related Work}
\label{sec:related}

\subsection{The Tsetlin Machine family}
In the Tsetlin Machine originally introduced by
Granmo~\cite{granmo2018tsetlin}, each class is represented by a set of
conjunctive clauses over Boolean literals, and the decision of the
machine is computed as a voting sum over these clauses. Each literal
is controlled by a dedicated Tsetlin automaton whose include/exclude
action is adjusted through a reinforcement procedure driven by two
forms of feedback. Three practical advantages have been consistently
reported over neural-network alternatives in constrained deployments:
rule-level interpretability, integer-only inference, and a training
dynamic that is inherently robust to the vanishing-gradient problem.

Several extensions have subsequently been published. A patch-based
feedback scheme was proposed for two-dimensional images in the
convolutional variant~\cite{granmo2019convolutional}. Sparse literal
budgets have been investigated in several forms:
Bhattarai~\emph{et~al.}~\cite{bhattarai2023concise} introduced a soft
clause-size constraint under which literals are expelled once a clause
grows beyond a target size, and drop-clause regularisation was
proposed for the Sparse Tsetlin
Machine~\cite{sharma2023drop}. For edge inference,
Maheshwari~\emph{et~al.}~\cite{maheshwari2023redress} developed
REDRESS, which combines include-encoding compression with a custom
inference kernel and has been reported to deliver up to $5700\times$
energy savings relative to binarised neural networks on an
STM32F746G-DISCO microcontroller. The most recent variant, on which
the present work builds, is the Fuzzy-Pattern Tsetlin Machine
(FPTM)~\cite{hnilov2025fuzzy}, which is described in detail in
Section~\ref{sec:fptm}.

\subsection{Intrusion-detection benchmarks}
The four public benchmarks used in this study have each been
established as de-facto standards for IoT and IoMT IDS research.
NSL-KDD was proposed by Tavallaee~\emph{et~al.}~\cite{tavallaee2009nslkdd}
as a cleaner, de-duplicated replacement for the KDD-CUP-99 corpus.
TON\_IoT~\cite{moustafa2021ton}, released by the University of New
South Wales, provides multi-protocol IoT traffic over ten attack
categories and is widely adopted in edge-IDS studies because of its
protocol heterogeneity. WUSTL-EHMS-2020~\cite{hady2020wustl} was
produced at Washington University in St.~Louis and is the first
public dataset to combine network-flow metrics with patient biomedical
readings; it is regarded as the most challenging of the four because
of its small size and severe class imbalance. MedSec-25 is a more
recent IoMT corpus that provides labelled flow records over five
attack categories and is included here to probe generalisation to a
newly proposed benchmark.

%% ========================================================================
\section{Binarisation for Tsetlin Machines}
\label{sec:binarisers}
\label{sec:glade}

Because every Tsetlin Machine variant operates on Boolean literals,
every continuous and categorical feature of a dataset must first be
mapped to one or more bits. Only a handful of binarisation schemes
are in wide use. The three that are encountered most often in the TM
literature are briefly reviewed below, after which the proposed GLADE
binariser is presented in detail.

\subsection{TMU \texttt{StandardBinarizer}}
\label{sec:bin-standard}
The reference binariser shipped with the Tsetlin Machine Unified
library~\cite{cair2024tmu}, denoted \texttt{StandardBinarizer} in the
reference code, implements the continuous-input scheme originally
proposed by Abeyrathna \emph{et~al.}~\cite{abeyrathna2019continuous} and
later refined into an adaptive sparse variant~\cite{abeyrathna2021adaptive}.
A threshold is placed at every distinct value observed in the training
column, subject to a user-specified cap
\texttt{max\_bits\_per\_feature}. The scheme is fully unsupervised, and
it has been adopted as the default encoder in several convolutional-TM
studies~\cite{granmo2019convolutional,bhattarai2023concise}. Two
limitations have been reported when the scheme is applied to IoT/IoMT
traffic data: the per-feature cap is shared across columns, so that
near-constant columns waste bits while high-cardinality columns are
under-resolved; and sparse columns in which zero dominates the mass
inherit the cap unchanged, which inflates the total bit count without
delivering a proportional gain in macro F1.

\subsection{K-bin quantile discretisation}
\label{sec:bin-kbins}
Equal-frequency quantile binning, available in scikit-learn as
\texttt{KBinsDiscretizer} with \texttt{strategy="quantile"}, places a
fixed number of thresholds at evenly spaced quantiles of each
training column and emits one bit per interior edge. The scheme is
label-free, compact, and has been adopted as the default encoder in
REDRESS~\cite{maheshwari2023redress} and in a number of subsequent
TM studies targeting tabular data. Its principal weakness is that the
quantile edges of sparse columns collapse onto the same value: on a
column in which $80\%$ of the samples are zero, every inner quantile
falls below the non-zero region, and a single effective threshold is
produced instead of the requested $n_\text{bins}$ thresholds. The
resulting bit budget is wasted, and the TM is required to compensate
through additional clauses.

\subsection{Thermometer encoding}
\label{sec:bin-thermo}
In a thermometer encoder, a fixed number of evenly spaced thresholds
are placed between the training minimum and maximum of each feature,
and the successive sign bits are concatenated so that bit $k$
indicates whether the feature exceeds the $k$-th threshold. Order
information is therefore preserved across bits, and the scheme has
been used in early TM applications because of its simplicity. The
practical limitation observed on IoT datasets is that the total bit
count grows linearly in both the number of features and in the
per-feature resolution: on the benchmarks of
Section~\ref{sec:results}, the thermometer encoder has been measured
to emit between $570$ and $1770$ bits, which is two to ten times the
output of the other schemes, so that the memory footprint of the
downstream TM is no longer competitive with tree-ensemble baselines.

\subsection{Proposed: the GLADE binariser}
\label{sec:glade-intro}
The proposed binariser, GLADE, is designed to remedy the three
limitations noted above simultaneously. Four design constraints are
imposed from the outset:
\begin{enumerate}
  \item \textbf{Unsupervised.} No class labels are read during the
        fitting procedure; every threshold is derived from the
        statistics of the training features only.
  \item \textbf{Single knob.} A single integer hyperparameter
        $n_\text{bins}$ is exposed to the user, and the same value is
        used for every feature of every dataset.
  \item \textbf{Adaptive per feature.} The number of bits actually
        allocated to a column is permitted to fall below $n_\text{bins}$
        when the column does not carry enough information to justify
        the full budget.
  \item \textbf{MCU-portable.} The fitted state is stored with
        \texttt{float32} thresholds and \texttt{int32} feature
        indices, so that a C reader on a 32-bit microcontroller can
        memory-map the file directly from flash.
\end{enumerate}
The single-knob requirement is enforced throughout the paper: for
every experiment and every dataset, $n_\text{bins}$ is set to $15$
and is never tuned. The next four subsections describe the four stages
of the binariser in turn: the adaptive bit budget
(Sec.~\ref{sec:glade-budget}), the conditional gap-aware threshold
placement (Sec.~\ref{sec:glade-gap}), the entropy-based dead-bit
elimination (Sec.~\ref{sec:glade-dead}), and the local gap
perturbation (Sec.~\ref{sec:glade-perturb}). The transform at
inference time and the TBF persistent-storage format are given in
Sec.~\ref{sec:glade-transform} and Sec.~\ref{sec:glade-tbf}.

\subsection{Adaptive bit budget}
\label{sec:glade-budget}
Let $c$ denote a feature column and let $u(c)$ be the number of
distinct values observed in the training sample of that column. For
columns with $u(c) \le n_\text{bins}$, the column is treated as a
categorical one-hot feature and exactly $u(c) - 1$ thresholds are
placed at the midpoints between successive unique values. For
columns with $u(c) > n_\text{bins}$, the budget is computed under a
\emph{hybrid} rule. Let $z(c) \in [0,1]$ denote the fraction of
training samples at which the column takes the value zero, and let
$u_{nz}(c)$ denote the number of distinct values among the non-zero
samples.

\emph{Case 1 --- dense continuous features} ($z(c) \le 0.3$ and
$u(c) > n_\text{bins}$). These columns typically hold genuine
continuous measurements such as packet rates, inter-arrival times or
biomedical readings. They receive the full
KBinsQuantile-equivalent budget
\begin{equation}
  b(c) \;=\; n_\text{bins},
  \label{eq:dense-budget}
\end{equation}
so that the per-feature resolution matches the strongest
quantile-based baseline. This case was added in the v2 revision of
the binariser after the v1 log-budget was observed to under-resolve
dense continuous features and to lose
0.02--0.03 macro F1 on the smaller WUSTL-EHMS-2020 dataset.

\emph{Case 2 --- sparse or near-categorical features}
($z(c) > 0.3$). These columns are typical of IoT traffic data, where
one-hot flag encodings and byte-count features concentrate a large
proportion of mass on zero. An effective unique count is defined as
\begin{equation}
  \mathrm{eff}(c) \;=\; \max \bigl(\, u_{nz}(c)\,(1 - z(c))^2,\; 2\bigr),
  \label{eq:effunique}
\end{equation}
and the bit budget is
\begin{equation}
  b(c) \;=\; \mathrm{clip}\Bigl(\lceil \log_2 \mathrm{eff}(c)\rceil + 2,\;
             1,\; n_\text{bins}\Bigr).
  \label{eq:budget}
\end{equation}
The term $(1 - z(c))^2$ penalises columns that are both sparse and
skewed: a feature that is zero 95\% of the time and takes two other
values rarely is assigned a budget close to $1$. The constant $+2$
term grants two additional bits to any sparse column so that the
zero versus non-zero distinction is always representable and so that
rare non-zero values remain discriminable.

The hybrid rule guarantees that GLADE allocates at least as many
bits as KBinsQuantile on every dense continuous feature, while
retaining the adaptive budget behaviour that is important on the
sparse one-hot flag columns of IoT datasets. This hybridisation is
the single most important change between v1 and v2 of the binariser
and is examined in the GLADE-component ablation of
Section~\ref{sec:ablation}.

\subsection{Conditional gap-aware threshold placement}
\label{sec:glade-gap}
Once $b(c)$ has been selected, $b(c)$ thresholds are placed along the
value range of the column. The default choice is evenly spaced
quantiles. This default can be biased when the data contain
structural gaps, that is, intervals of the value axis that are empty
or nearly empty due to protocol boundaries, sensor resolution or
categorical substructure hidden inside a nominally continuous column.
On such features, a threshold that falls inside a structural gap is
wasteful because it produces the same bit for every sample on either
side of the gap. Conversely, on smooth continuous features a naive
snap to the nearest gap introduces floating-point bias and leaves the
thresholds clustered unevenly.

A conditional nudge is therefore used. Let $U(c) = \{u_1,\dots,u_k\}$
be the sorted vector of unique non-zero values and let
$g_i = u_{i+1} - u_i$ be the vector of successive gaps. The column is
declared to possess structural gaps when
\begin{equation}
  \max_i g_i \;>\; 5 \cdot \mathrm{median}_i\,g_i,
  \label{eq:structural}
\end{equation}
and only in that case are the candidate quantile edges snapped to
the midpoints $(u_i + u_{i+1})/2$ of the largest gaps that lie in
their neighbourhood. Features for which
Eq.~\ref{eq:structural} fails retain their quantile edges unchanged.
The factor of five is selected empirically and is kept fixed across
all datasets in Section~\ref{sec:results}.

\subsection{Dead-bit elimination}
\label{sec:glade-dead}
A further post-processing step is applied after the thresholds have
been placed. In v1 of the binariser, a threshold was considered
dead when fewer than $1\%$ of the training samples lay on either
side of it. The v2 revision replaces this activation-fraction rule
with an \emph{entropy} criterion computed directly on the resulting
Boolean bit. For each candidate threshold $t_j$ the resulting bit
is $b_j(\mathbf{x}) = [x_{f_j} \ge t_j]$ and its Bernoulli entropy
is
\begin{equation}
  H(b_j) \;=\; -p\,\log p \;-\; (1-p)\log(1-p),
  \quad p = \Pr_{\mathbf{x}\in X_{\text{tr}}}[b_j(\mathbf{x})=1].
  \label{eq:bit-entropy}
\end{equation}
A threshold is discarded when $H(b_j) < 0.05$ nats. Dead thresholds
are replaced by the midpoint of the widest surviving interval,
iteratively, until either the full budget $b(c)$ is reached again or
no wider interval can be split.

The entropy criterion is strictly more appropriate for imbalanced
datasets than the activation-fraction rule. On WUSTL-EHMS-2020 the
smallest class represents only $5.6\%$ of the training samples, and
a threshold whose Boolean bit activates on just $0.7\%$ of samples
may still be the single discriminator between that class and the
normal traffic. Rule~(\ref{eq:bit-entropy}) preserves such bits
because it keeps any threshold whose bit has non-trivial Bernoulli
entropy, not only those that activate on more than $1\%$ of samples.

\subsection{Local gap perturbation}
\label{sec:glade-perturb}
A further refinement is applied to the quantile edges. Let
$t_1 < t_2 < \cdots < t_{b(c)}$ denote the sorted thresholds
returned by the previous stages, and let $\ell_i = t_i - t_{i-1}$
be the local half-gap to the left, with $\ell_1$ extending to the
training minimum and a symmetric definition for the right half-gap.
For each threshold $t_i$ three candidate placements are scored:
\begin{equation}
  \mathcal{C}_i \;=\; \bigl\{\,
     t_i - \tfrac{\ell_i}{4},\;\;
     t_i,\;\;
     t_i + \tfrac{r_i}{4}\,
  \bigr\},
  \label{eq:perturb}
\end{equation}
and the placement that maximises the Bernoulli variance
$p(1-p)$ of the resulting bit is selected. The perturbation is
deterministic and runs in $\mathcal{O}(b(c))$ time per column. Its
effect is to move quantile edges into the high-density neighbourhood
of the feature distribution, away from the rigid equal-frequency
placement that KBinsQuantile would use.

\subsection{Transform and complexity}
\label{sec:glade-transform}
At inference time, the fitted binariser is stored as two flat arrays:
a feature-index array $\mathbf{f} \in \mathbb{Z}^{N_\text{bits}}$ and
a threshold array $\mathbf{t} \in \mathbb{R}^{N_\text{bits}}$, where
$N_\text{bits} = \sum_c b(c)$ is the total number of output bits.
For a new sample $\mathbf{x}$, the $i$-th output bit is computed as
\begin{equation}
  \mathrm{bit}_i(\mathbf{x}) \;=\;
      \bigl[\, x_{f_i} \,\ge\, t_i \,\bigr],
  \label{eq:transform}
\end{equation}
which is a single fused comparison per bit, without branches and
without per-feature bookkeeping. The asymptotic cost of the
transformation is $\mathcal{O}(N_\text{bits})$ per sample and the
total memory footprint of the binariser is
$4 \cdot N_\text{bits}$ bytes for the thresholds plus
$2 \cdot N_\text{bits}$ bytes for the feature indices, which is at
most $2$~KB for any of the datasets used in Section~\ref{sec:results}.

%% ========================================================================
\section{The Fuzzy-Pattern Tsetlin Machine}
\label{sec:fptm}

The downstream classifier used in this paper is the Fuzzy-Pattern
Tsetlin Machine (FPTM) of Hnilov~\cite{hnilov2025fuzzy}. It is used
verbatim, through the public reference implementation, and no aspect
of its training procedure is modified. A self-contained description
is given below so that the remainder of the paper can be followed
without consulting external sources; the exposition follows the
visual style of REDRESS~\cite{maheshwari2023redress}.

\begin{figure*}[t]
\centering
\IfFileExists{fig/fig_fptm.pdf}
  {\includegraphics[width=0.95\textwidth]{fig/fig_fptm.pdf}}%
  {\fbox{\parbox{0.9\textwidth}{\vspace*{60pt}%
    \centering\textbf{[Fig. \ref{fig:fptm}]}\\[4pt]%
    \emph{Export \texttt{fig\_fptm.drawio} to
    \texttt{fig/fig\_fptm.pdf} to replace this placeholder.}%
    \vspace*{60pt}}}}%
\caption{Fuzzy-Pattern Tsetlin Machine inference pipeline (editable
         source in \texttt{fig\_fptm.drawio}). Raw features
         $\mathbf{x}$ are booleanised by GLADE into literal bits and
         their complements. For every clause, the number of mismatched
         literals $m$ is counted by a single bit-population count over
         the packed literal words, and the graded clause output
         $o = \max(\kappa - m,\, 0)$ is obtained through a clamped
         subtraction. The multiclass voting block illustrates the
         bit-parallel evaluation of the $2CK$ clauses of a $K$-class
         model: each class holds $C$ positive-polarity and $C$
         negative-polarity clauses whose outputs are signed and summed
         into a class vote $V_k$, and the predicted class is the
         $\arg\max$ over $V_k$.}
\label{fig:fptm}
\end{figure*}

\subsection{Fuzzy clause evaluation}
\label{sec:fptm-clause}
An FPTM maintains, for every output class $k$ and for every polarity
$\sigma \in \{+,-\}$, a fixed number $C$ of \emph{clauses}. A clause
$\phi$ is a set of \emph{literals} drawn from the Boolean input
$\mathbf{x} \in \{0,1\}^{N}$ and its complement $\neg\mathbf{x}$. The
inclusion state of each literal is controlled by a dedicated Tsetlin
automaton whose action is \emph{include} when the internal state has
reached a learned threshold and \emph{exclude} otherwise. Let
$I_\phi \subseteq \{1,\dots,N\}$ denote the set of positive literals
actually included in $\phi$ and let $E_\phi$ denote the set of
negated literals included in $\phi$. The \emph{mismatch count} of
$\phi$ on the input $\mathbf{x}$ is then defined as
\begin{equation}
  m(\phi, \mathbf{x}) \;=\;
    \bigl|\{\, i \in I_\phi : x_i = 0 \,\}\bigr|
    \;+\; \bigl|\{\, i \in E_\phi : x_i = 1 \,\}\bigr|,
  \label{eq:fptm-mismatch}
\end{equation}
and the \emph{graded clause output} is
\begin{equation}
  o(\phi, \mathbf{x}) \;=\;
    \max\bigl(\, \kappa_\phi - m(\phi,\mathbf{x}),\;\; 0 \,\bigr),
  \label{eq:fptm-output}
\end{equation}
where the per-clause saturation cap is
$\kappa_\phi = \min\bigl(|I_\phi \cup E_\phi|,\; L_f\bigr)$ and
$L_f$ is a user-defined clamp. A clause therefore contributes to its
class vote even when a small number of its literals fail, with the
contribution decaying linearly in the mismatch count and saturating
at $\kappa_\phi$ when every included literal is matched. The smooth
degradation introduced by Eq.~\eqref{eq:fptm-output} is what enables
FPTM to reach the accuracy of earlier Tsetlin variants with between
$50$ and $400$ times fewer clauses~\cite{hnilov2025fuzzy}, and is the
reason that the tuned GLADE+FPTM configurations of
Section~\ref{sec:results} fit in $17$--$68$~kilobytes of on-disk
storage.

\subsection{Bit-parallel evaluation}
\label{sec:fptm-parallel}
A decisive property of Eq.~\eqref{eq:fptm-mismatch} is that it can be
computed as a single bit-population count over 64-bit packed words.
Let $\mathbf{w}^{I}_{\phi}$ and $\mathbf{w}^{E}_{\phi}$ denote the
packed inclusion masks of the positive and negated literals of
$\phi$, and let $\mathbf{w}_{x}$ denote the packed input word. The
mismatch count is then
\begin{equation}
  m(\phi,\mathbf{x})
  \;=\; \operatorname{popcount}\!\bigl(\mathbf{w}^{I}_{\phi} \wedge \neg\mathbf{w}_{x}\bigr)
  \;+\; \operatorname{popcount}\!\bigl(\mathbf{w}^{E}_{\phi} \wedge \mathbf{w}_{x}\bigr),
  \label{eq:popcount}
\end{equation}
which requires a number of 64-bit operations proportional to
$\lceil N/64 \rceil$, a value that falls between $2$ and $6$ on the
datasets of Section~\ref{sec:results}. Because every clause accesses
the input word through read-only bitmasks, the $2CK$ clauses of a
$K$-class model are evaluated independently and can be dispatched in
parallel across CPU cores, across SIMD lanes of a vector-extension
(e.g. NEON, AVX2, SVE), or across the tiled processing elements of an
accelerator. The reference implementation exploits this parallelism
through a lightweight thread pool that assigns disjoint blocks of
clauses to worker threads, and the evaluation of a class vote reduces
to a fused signed accumulation of the clause outputs,
\begin{equation}
  V_k(\mathbf{x}) \;=\;
    \sum_{\phi \in \Phi_k^{+}} o(\phi,\mathbf{x})
    \;-\; \sum_{\phi \in \Phi_k^{-}} o(\phi,\mathbf{x}).
  \label{eq:classvote}
\end{equation}
The predicted class is $\hat{k} = \arg\max_{k} V_k(\mathbf{x})$. On
the single-core Raspberry~Pi~5 used in Section~\ref{sec:results}, the
per-sample inference latency of the GLADE+FPTM Numba JIT pipeline is
measured between $4.10$ and $7.42$~microseconds, while a Numba-compiled
DecisionTree baseline reaches $1.37$--$1.39$~microseconds, confirming that
the bit-parallel evaluation is the factor that makes the pipeline attractive
for deployment on constrained edge CPUs.

\subsection{Training feedback and role in the pipeline}
\label{sec:fptm-training}
At training time, $V_k$ is driven toward a target magnitude $T$
through feedback on the Tsetlin automata of every clause: clauses of
the target class are updated through Type~I feedback and clauses of
every non-target class are updated through Type~II feedback, both
gated by a specificity parameter $s$ and a literal cap $L$. The
feedback rule is given in Fig.~\ref{fig:fptm-feedback}, which is
transcribed directly from the public reference implementation; it is
not modified by the present work. The FPTM therefore enters the
pipeline as a drop-in classifier whose hyperparameters
$(C, T, s, L, L_f)$ are either held fixed across datasets in the
shared-baseline row or swept jointly in the tuned row under the
constraint $C \le 100$.

\begin{figure*}[t]
\centering
\IfFileExists{fig/fig_fptm_feedback.pdf}
  {\includegraphics[width=0.95\textwidth]{fig/fig_fptm_feedback.pdf}}%
  {\fbox{\parbox{0.9\textwidth}{\vspace*{60pt}%
    \centering\textbf{[Fig. \ref{fig:fptm-feedback}]}\\[4pt]%
    \emph{Export \texttt{fig\_fptm\_feedback.drawio} to
    \texttt{fig/fig\_fptm\_feedback.pdf} to replace this placeholder.}%
    \vspace*{60pt}}}}%
\caption{Fuzzy-Pattern feedback procedure. For every clause, the
         graded output $o \in [0,\kappa]$ is first computed; if
         $o>0$ the clause is reinforced (Type~Ia) through deterministic
         per-bit updates on the positive automaton $c[i]$ and the
         negated automaton $ci[i]$, and mismatched excluded literals
         are forgotten; if $o=0$ the clause receives $s$ random
         decrement steps on both automata (Type~Ib). Type~II feedback
         is applied for non-target classes. The figure is transcribed
         directly from the public reference implementation and is
         shown here for completeness; it is not modified by the
         present work.}
\label{fig:fptm-feedback}
\end{figure*}

\subsection{Complexity and learning-theoretic view}
\label{sec:fptm-theory}
A compact analysis of the pipeline is given here to complement the
empirical results of Section~\ref{sec:results}.

\emph{Time and memory complexity.} Let $N$ denote the total number
of GLADE bits, let $C$ denote the number of clauses per class per
polarity, and let $K$ denote the number of output classes. At
inference time, every clause is evaluated through
$\lceil N/64 \rceil$ packed logical operations followed by a
constant-time \texttt{popcount}, so that the per-sample inference
cost is
$
  \mathcal{O}\!\bigl(\,C\,K\,\lceil N/64 \rceil\,\bigr)
$
bit-parallel operations. For the values of $(C,K,N)$ used in
Section~\ref{sec:results}, this evaluates to between $5{,}000$ and
$40{,}000$ 64-bit operations per sample, which is three to five
orders of magnitude below the per-sample cost of a tree-ensemble
baseline of comparable accuracy. The memory footprint of the
trained model is bounded above by $2CK\,\lceil N/64 \rceil$ bytes
for the packed literal bitmasks plus $\mathcal{O}(N)$ bytes for the
fitted GLADE state; the constants are small enough that the model
fits in at most $56$~kilobytes on every dataset studied, without any
post-training quantisation.

\emph{GLADE as an entropy-aware feature map.} The dead-bit
elimination step of Sec.~\ref{sec:glade-dead} discards any
threshold whose induced Boolean bit has Bernoulli entropy below
$0.05$~nats, which corresponds to a per-bit activation probability
$p$ outside the interval $(0.014,\,0.986)$. A bit that survives
pruning therefore activates on neither too few nor too many
training samples, and any budget that is not consumed by surviving
thresholds is refilled with midpoint thresholds on the widest
remaining intervals. The practical consequence is that the
majority of the literals seen by the downstream FPTM are
individually informative: no single literal is constant on the
training set, and very few literals fire on a pathologically
small or pathologically large fraction of samples.

\emph{Why the graded mismatch rule improves generalisation.}
Eq.~\eqref{eq:fptm-output} is a piecewise-linear function of the
mismatch count whose slope is $-1$ and whose ceiling is the
per-clause cap $\kappa_\phi$. A single literal failure therefore
decreases the clause contribution by exactly one, rather than
collapsing it to zero as in an all-or-nothing rule, and the total
change induced in the class vote $V_k(\mathbf{x})$ by a Hamming-distance
perturbation of size $\delta$ on the literal vector is bounded
above by $C\delta$, where $C$ is the number of clauses per class.
A single spurious literal therefore cannot disable a clause and
cannot flip a class decision that the remaining literals
unambiguously support. The same bounded-change property forces the
feedback procedure of Fig.~\ref{fig:fptm-feedback} to make smaller
incremental changes to the clauses, which is observed to reduce
the variance of the training dynamics across seeds (the five-seed
standard deviations reported in Sec.~\ref{sec:results} stay below
$0.011$ on every dataset). The $50$--$400\times$ clause-count
reduction reported by Hnilov~\cite{hnilov2025fuzzy} is interpreted
in this paper as a direct consequence of that bounded-sensitivity
property: because each clause tolerates a small number of literal
failures without silencing, fewer clauses are required to cover
the input space, and the compact model sizes of
Section~\ref{sec:results} are consequently a structural property
of the learner rather than an engineering artefact.

%% ========================================================================
\section{Experimental Setup}
\label{sec:setup}

\subsection{Datasets and preprocessing}
\label{sec:datasets}
The four public benchmarks used in this study are summarised in a
single combined view in Table~\ref{tab:datasets}. The same
preprocessing pipeline is applied to every dataset and to every
classifier so that no baseline is advantaged by a custom
feature-engineering step: a stratified $80/20$ train--test split is
drawn with a fixed seed of $42$; infinities are replaced with NaNs,
and NaNs are filled with the per-column median computed on the
training partition only; constant columns are removed on the basis
of the training-set standard deviation; and any column whose Pearson
correlation coefficient with an earlier kept column exceeds $0.99$
is discarded, with the correlation again computed on the training
partition. The same pipeline is shared between the Tsetlin path and
the non-Tsetlin baselines; the only difference is that the Tsetlin
path additionally passes its preprocessed features through GLADE,
whereas the other classifiers consume the raw (or standard-scaled)
numerical features directly. The last three columns of
Table~\ref{tab:datasets} report the final training and test sizes
after this pipeline, and the number of GLADE bits emitted with
$n_\text{bins}=15$.

\begin{table*}[t]
\caption{Combined view of the four public benchmarks used in this
         study. Record and feature counts are reported before
         preprocessing. The last three columns are reported after the
         shared preprocessing pipeline (median fill, constant-column
         removal, and high-correlation dedup). The ``GLADE bits''
         column is the total output width of the fitted GLADE
         binariser at $n_\text{bins}=15$, which is the input width
         seen by the downstream FPTM.}
\label{tab:datasets}
\centering\small
\begin{tabular}{lrrrrrrr}
\toprule
\textbf{Dataset} & \textbf{Records} & \textbf{Raw feat.} &
\textbf{Classes} & \textbf{Domain} &
\textbf{Train} & \textbf{Test} & \textbf{GLADE bits} \\
\midrule
NSL-KDD~\cite{tavallaee2009nslkdd}       & $148{,}517$ & $41$ & $5$  & IoT network flow             & $118{,}813$ & $29{,}704$  & $232$ \\
TON\_IoT~\cite{moustafa2021ton}          & $211{,}043$ & $44$ & $10$ & IoT multi-protocol           & $168{,}834$ & $42{,}209$  & $115$ \\
MedSec-25                                & $554{,}534$ & $78$ & $5$  & IoMT flow                    & $443{,}627$ & $110{,}907$ & $349$ \\
WUSTL-EHMS-2020~\cite{hady2020wustl}     & $16{,}318$  & $44$ & $3$  & IoMT network $+$ biomedical  & $13{,}054$  & $3{,}264$   & $222$ \\
\bottomrule
\end{tabular}
\end{table*}

\subsection{FPTM configuration}
\label{sec:fptm-config}
Two FPTM configurations are evaluated. The \emph{shared-baseline}
configuration reuses a single hyperparameter set across every
dataset ($C{=}80$, $T{=}10$, $s{=}61$, $L{=}60$, $L_f{=}10$,
$E{=}30$~epochs), where $C$ is the number of clauses per class per
polarity, $T$ the voting target, $s$ the specificity, $L$ the
literal cap, $L_f$ the fuzzy clamp, and $E$ the number of training
epochs. The \emph{tuned} configuration is obtained by a random
joint search of $500$ configurations per dataset under the
constraint $C\le 100$ with $C$ even, parallelised across $96$ CPU
workers with a fixed seed of $42$; the best macro-F1 configurations
returned by the search are
$(C,T,s,L,L_f,E)=(90,8,100,30,5,80)$ on WUSTL-EHMS-2020,
$(78,8,50,30,8,200)$ on NSL-KDD,
$(78,3,50,30,5,300)$ on TON\_IoT, and
$(96,15,150,20,5,80)$ on MedSec-25.

\subsection{Model storage: the Template Binary Format}
\label{sec:tbf}
Both the fitted GLADE state and the trained FPTM clauses are
serialised to disk through a single self-describing binary format,
referred to here as the Template Binary Format (TBF). The format is
used as the sole on-disk representation of the whole model,
comprising the thresholds and feature indices of the binariser, the
packed literal bitmasks of the trained clauses, the automaton states
at the include limit, and every scalar hyperparameter required at
inference time. A one-byte tag is placed before every value to
denote one of the primitive types (booleans, signed and unsigned
integers of $8$--$64$ bits, single- and double-precision floats,
UTF-8 strings, byte strings, lists, string-keyed dictionaries, and
multi-dimensional \texttt{ndarrays}). Every length field is encoded
as an unsigned variable-length integer (LEB128), and homogeneous
integer lists — which arise naturally in the clause representation
as lists of literal indices — are automatically promoted to the
narrowest unsigned integer \texttt{ndarray} dtype that can hold
their maximum value. The end-to-end effect is a reduction of the
serialised model size by $40$--$60\%$ relative to a JSON baseline
without any change to the inference semantics. Every byte order is
little-endian and no platform-specific padding is emitted, so that
the same file that is produced by the reference Python writer on a
training server is read directly by a portable C deserialiser of
approximately $150$ lines on an ARM microcontroller.

\subsection{Baselines}
Twelve baselines are trained under the same preprocessing pipeline
and the same single-core protocol: XGBoost (small and full),
RandomForest (small and full), LogisticRegression, LinearSVM,
DecisionTree, GaussianNB, k-NN with $k=5$, and MLP (tiny, small,
med). The XGBoost configurations are the default multi-class loss,
$25$ and $100$ trees respectively, and the RandomForest
configurations use $25$ and $100$ trees with depths $6$ and $12$. The
linear models use the default L2 regularisation and the solver
LBFGS with a $300$-iteration cap. All baselines are constructed
through scikit-learn~\cite{pedregosa2011scikit} or the official
XGBoost Python package.

\subsection{Evaluation protocol}
Every training and inference run is constrained to a single CPU core
by exporting \texttt{OPENBLAS\_NUM\_THREADS=1},
\texttt{MKL\_NUM\_THREADS=1}, \texttt{OMP\_NUM\_THREADS=1} and
\texttt{NUMEXPR\_NUM\_THREADS=1} before the Python process imports
NumPy, by passing \texttt{n\_jobs=1} to every scikit-learn and
XGBoost constructor, and by setting \texttt{NUMBA\_NUM\_THREADS=1}
for the Numba JIT FPTM runs. Wall-clock times are measured with
\texttt{time.perf\_counter}. Macro F1 is computed with scikit-learn's
\texttt{f1\_score(average="macro")}; the per-class reports and the
weighted F1 are also recorded but are not shown in the tables for
reasons of space. The reported timings separate the binarisation
phase from the TM training phase so that fair comparisons to sklearn
\texttt{fit} (which includes preprocessing) can be computed as
\texttt{binarize\_time + tm\_fit\_time}.


% In your preamble: \usepackage{booktabs}
% In your preamble:


\begin{figure*}[t]
\centering
\begin{subfigure}[b]{0.48\textwidth}
    \centering
    \includegraphics[width=\linewidth]{boxplot_f1_glade_v2_NSL-KDD.png}
    \caption{NSL-KDD}
    \label{fig:box_nslkdd}
\end{subfigure}
\hfill
\begin{subfigure}[b]{0.48\textwidth}
    \centering
    \includegraphics[width=\linewidth]{boxplot_f1_glade_v2_TON_IoT.png}
    \caption{TON-IoT}
    \label{fig:box_toniot}
\end{subfigure}

\vspace{0.5em}

\begin{subfigure}[b]{0.48\textwidth}
    \centering
    \includegraphics[width=\linewidth]{boxplot_f1_glade_v2_MedSec-25.png}
    \caption{MedSec-25}
    \label{fig:box_medsec}
\end{subfigure}
\hfill
\begin{subfigure}[b]{0.48\textwidth}
    \centering
    \includegraphics[width=\linewidth]{boxplot_f1_glade_v2_WUSTL.png}
    \caption{WUSTL-EHMS-2020}
    \label{fig:box_wustl}
\end{subfigure}

\caption{Macro-F1 vs.\ number of bins ($n_{\text{bins}}$) across all four datasets. GLADE\_v2 (gray boxes, 5-fold cross-validation) is compared against StandardBin (TMU, orange $\times$) and KBinsUniform (green $\bullet$) baselines. GLADE\_v2 converges faster (fewer bins) and consistently matches or outperforms both baselines, with the largest margins on TON-IoT and MedSec-25.}
\label{fig:binarizer_boxplots}
\end{figure*}
\begin{table}[t]
\centering
\caption{Model configurations. Starred ($^\ast$) variants use hyperparameters matched to the TM's clause count $C$ (or literal budget $L$) for a fair comparison.}
\label{tab:model_configs}
\begin{tabular}{ll}
\toprule
Model & Configuration \\
\midrule
\multicolumn{2}{l}{\textit{ML baselines (library defaults)}} \\
\midrule
XGBoost              & $n{=}100$, depth${=}6$, lr${=}0.1$ \\
RandomForest         & $n{=}100$, depth${=}$None \\
DecisionTree         & depth${=}$None, gini \\
MLP\_tiny            & hidden${=}(32)$, ReLU, 200 ep. \\
MLP\_small           & hidden${=}(64)$, ReLU, 200 ep. \\
MLP\_med             & hidden${=}(128,64)$, ReLU, 200 ep. \\
kNN\_5               & $k{=}5$, uniform, Minkowski \\
LogisticRegression   & L2, $C{=}1.0$, lbfgs \\
LinearSVM            & $C{=}1.0$, squared hinge \\
NaiveBayes           & GaussianNB \\
\midrule
\multicolumn{2}{l}{\textit{TM-matched variants}} \\
\midrule
XGBoost\_Cmatched$^\ast$      & $n_{\text{est}} = C$ \\
RandomForest\_Cmatched$^\ast$ & $n_{\text{est}} = C$ \\
MLP\_C$^\ast$                 & hidden ${=}(C)$ \\
MLP\_2C$^\ast$                & hidden ${=}(2C, C)$ \\
DecisionTree\_Lmatched$^\ast$ & max depth ${=}L$ \\
\bottomrule
\end{tabular}
\end{table}

\begin{table}[t]
\centering
\caption{Tsetlin Machine hyperparameters per dataset. $C$: clauses per class, $T$: voting threshold, $S$: specificity, $L$: max literals per clause, $LF$: literal forgetting, $E$: epochs.}
\label{tab:tm_hparams}
\begin{tabular}{lcccccc}
\toprule
Dataset & $C$ & $T$ & $S$ & $L$ & $LF$ & $E$ \\
\midrule
NSL-KDD         & 90  & 12 & 200 & 40 & 8  & 80  \\
TON-IoT         & 100 & 15 & 20  & 25 & 8  & 200 \\
MedSec-25       & 80  & 12 & 75  & 30 & 8  & 300 \\
WUSTL-EHMS-2020 & 60  & 8  & 300 & 50 & 15 & 200 \\
\bottomrule
\end{tabular}
\end{table}


% In your preamble: \usepackage{booktabs}

\begin{table*}[t]
\centering
\caption{Per-dataset comparison of binarizers at $n_{\text{bins}}{=}15$. GLADE\_v2 uses roughly half the bits of the TMU StandardBin while matching or exceeding its accuracy and macro-F1. Best value per metric per dataset in \textbf{bold}.}
\label{tab:binarizer_comparison}
\begin{tabular}{l ccc ccc ccc}
\toprule
 & \multicolumn{3}{c}{GLADE\_v2} & \multicolumn{3}{c}{StandardBin (TMU)} & \multicolumn{3}{c}{KBinsUniform} \\
\cmidrule(lr){2-4} \cmidrule(lr){5-7} \cmidrule(lr){8-10}
Dataset & Bits & Acc & F1 & Bits & Acc & F1 & Bits & Acc & F1 \\
\midrule
NSL-KDD         & \textbf{254} & \textbf{0.9955} & 0.9311          & 472 & 0.9941 & \textbf{0.9342} & 448 & 0.9854 & 0.8690 \\
TON-IoT         & \textbf{120} & \textbf{0.9688} & \textbf{0.9385} & 186 & 0.8416 & 0.7947          & 177 & 0.6411 & 0.5663 \\
MedSec-25       & \textbf{349} & \textbf{0.9703} & \textbf{0.8681} & 698 & 0.8830 & 0.7265          & 653 & 0.8461 & 0.5532 \\
WUSTL-EHMS-2020 & \textbf{266} & \textbf{0.9393} & \textbf{0.7630} & 313 & 0.9381 & 0.7388          & 296 & 0.9130 & 0.6259 \\
\bottomrule
\end{tabular}
\end{table*}
%% ========================================================================
\section{Results}
\label{sec:results}

\subsection{Per-sample inference latency}
\label{sec:res-latency}

Table~\ref{tab:full_latency} reports the per-sample inference latency
of every model on the Raspberry~Pi~5, measured under the single-core
protocol of Section~\ref{sec:setup}. Each model is evaluated in its
native deployment form: pure NumPy for the ML baselines and the Python
FPTM variant, and Numba JIT for the FPTM Numba variant. Three
observations are drawn.

First, DecisionTree compiled with Numba achieves 1.37--1.39~$\mu$s,
the lowest-latency entry by a large margin. Second, the GLADE+FPTM
Numba kernel achieves 4.10--7.42~$\mu$s, which is 4--6$\times$ slower
than DT Numba but maintains higher macro-F1 on three of four datasets.
Third, the pure-Python FPTM remains competitive with the NumPy baselines
but is 11--17$\times$ slower than the Numba variant.

\begin{table*}[t]
\centering\small
\caption{Per-sample inference latency ($\mu$s) on a single Cortex-A76 core of
         the Raspberry~Pi~5 at 2.4~GHz, measured under the single-core
         protocol (100-iteration warm-up, timed loop), using pure-numpy
         for the ML baselines. Best per-dataset value in \textbf{bold}.}
\label{tab:full_latency}
\begin{tabular}{lrrrr}
\toprule
\textbf{Model} & \textbf{TON\_IoT} & \textbf{MedSec-25} & \textbf{WUSTL} & \textbf{NSL-KDD} \\
\midrule
RandomForest                 & $847.97$    & $891.24$    & $840.31$   & $887.45$   \\
RandomForest\_Cmatched       & $422.66$    & $417.14$    & $405.28$   & $486.10$   \\
DecisionTree                 & $6.96$      & $6.63$      & $7.16$     & $8.29$     \\
DecisionTree\_Lmatched       & $8.60$      & $8.91$      & $10.96$    & $15.06$    \\
LogisticRegression           & $7.86$      & $8.06$      & $7.25$     & $7.94$     \\
LinearSVM                    & $7.86$      & $8.22$      & $7.31$     & $8.05$     \\
NaiveBayes                   & $24.30$     & $23.96$     & $20.54$    & $26.91$    \\
k-NN ($k{=}5$)               & $21675.57$  & $96872.75$  & $1017.69$  & $45862.43$ \\
MLP\_tiny                    & $12.68$     & $13.17$     & $12.18$    & $14.34$    \\
MLP\_small                   & $18.81$     & $19.45$     & $18.10$    & $21.16$    \\
MLP\_med                     & $20.90$     & $21.85$     & $20.27$    & $25.07$    \\
MLP\_C                       & $15.48$     & $16.53$     & $14.27$    & $21.26$    \\
MLP\_2C                      & $18.23$     & $21.56$     & $17.14$    & $30.12$    \\
\midrule
\textbf{DecisionTree (Numba)}        & $\mathbf{1.383}$ & $\mathbf{1.367}$ & $\mathbf{1.374}$ & $\mathbf{1.393}$ \\
\midrule
\textbf{GLADE+FPTM (Python)} & $88.52$     & $69.67$     & $51.79$    & $67.35$    \\
\textbf{GLADE+FPTM (Numba)}  & $6.722$     & $7.415$     & $4.101$    & $6.314$    \\
\bottomrule
\end{tabular}
\end{table*}

\begin{table*}[t]
\centering\small
\caption{Per-sample latency on a single Cortex-A76 core of the
         Raspberry~Pi~5 (2.4~GHz), single-threaded, 500-sample warm-up,
         3\,000-sample timed window. Two runtime implementations of the
         identical FPTM algorithm plus the fastest pure-numpy baseline.
         \emph{Python}: optimised FBZEngine (zero-alloc buffers, flat clause
         matrix); \emph{Numba}: \texttt{@njit(cache=True)} with SWAR popcount.
         Speed-up relative to Python numpy.}
\label{tab:pi5_runtime}
\begin{tabular}{lrrrrrrrr}
\toprule
 & \multicolumn{2}{c}{\textbf{WUSTL}} & \multicolumn{2}{c}{\textbf{NSL-KDD}} & \multicolumn{2}{c}{\textbf{TON\_IoT}} & \multicolumn{2}{c}{\textbf{MedSec-25}} \\
\cmidrule(lr){2-3}\cmidrule(lr){4-5}\cmidrule(lr){6-7}\cmidrule(lr){8-9}
\textbf{Runtime} & $\mu$s & speed-up & $\mu$s & speed-up & $\mu$s & speed-up & $\mu$s & speed-up \\
\midrule
Python numpy (FBZEngine) & $39.40$ & $1.0\times$ & $51.30$ & $1.0\times$ & $67.40$ & $1.0\times$ & $53.70$ & $1.0\times$ \\
Numba JIT (SWAR)         & $4.101$ & $9.6\times$ & $6.314$ & $8.1\times$ & $6.722$ & $10.0\times$ & $7.415$ & $7.2\times$ \\
\midrule
\textbf{DecisionTree (Numba)} & $\mathbf{1.374}$ & $\mathbf{28.7\times}$ & $\mathbf{1.393}$ & $\mathbf{36.8\times}$ & $\mathbf{1.383}$ & $\mathbf{48.7\times}$ & $\mathbf{1.367}$ & $\mathbf{39.3\times}$ \\
\bottomrule
\end{tabular}
\end{table*}

\subsection{Why compilation unlocks the hardware}
\label{sec:res-pynumba}

The same FPTM inference algorithm is implemented twice---once through a
NumPy-based Python function, once through a Numba JIT kernel---to isolate
the effect of the runtime. Both produce bit-identical predictions on the
full test partition of every dataset.

\begin{table}[t]
\centering\small
\caption{Same-algorithm, same-model, same-data runtime comparison on
         the Pi~5 (Cortex-A76). Speed-up is relative to the Python
         numpy baseline. All implementations produce bit-identical predictions.}
\label{tab:pyvsnumba}
\begin{tabular}{llrr}
\toprule
\textbf{Platform} & \textbf{Runtime} & \textbf{WUSTL ($\mu$s)} & \textbf{Speed-up} \\
\midrule
\multirow{2}{*}{Pi~5 (A76)}
 & Python numpy & $51.79$ & $1.0\times$ \\
 & Numba JIT    & $4.101$ & $\mathbf{12.6\times}$ \\
\bottomrule
\end{tabular}
\end{table}

Table~\ref{tab:pyvsnumba} reports a 9.6$\times$ to 12.6$\times$ speed-up
from the Numba JIT kernel over the pure-Python numpy implementation. Because
the algorithm is identical, the speed-up is entirely attributable to the
runtime: removal of per-call interpreter dispatch (2--5~$\mu$s per NumPy
call on the Cortex-A76), elimination of per-sample heap allocation, and
L1 cache residency of the clause bank after JIT compilation.

\subsection{Single-core throughput}
\label{sec:res-throughput}

Table~\ref{tab:throughput} reports the sustained single-core throughput of
the GLADE+FPTM Numba kernel on the Pi~5 and the projected latency on the
Pi~Zero~2W (Cortex-A53, 1.0~GHz). The Pi~5 single-core throughput is
between 134{,}864 and 243{,}843 predictions per second; across all four
cores, the Pi~5 sustains 539{,}000 to 975{,}000 predictions per second.
The projected Pi~Zero~2W latencies (16.4--29.7~$\mu$s) are well below the
minimum Ethernet frame transit time at 100~Mbit/s (0.51~ms), confirming
that the pipeline is viable for intrusion detection on access-link hardware.

\begin{table*}[t]
\centering\small
\caption{Sustained single-core throughput on Raspberry Pi~5 (measured,
         Numba JIT) and projected per-sample latency and throughput on the
         Raspberry Pi~Zero~2W (Cortex-A53, 1.0~GHz). Projected values are
         obtained by scaling Pi~5 latency by a factor of $4.0\times$
         (combined clock ratio $2.4\times$ and IPC ratio $1.67\times$,
         A76 vs A53). Four-core Pi~5 throughput assumes linear scaling.}
\label{tab:throughput}
\begin{tabular}{lrrrr}
\toprule
\textbf{Dataset}
  & \textbf{Pi~5 single core}
  & \textbf{Pi~5 four cores (proj.)}
  & \textbf{Pi~Zero~2W latency (proj.)}
  & \textbf{Pi~Zero~2W single core (proj.)} \\
\midrule
NSL-KDD         & $158{,}503~\text{s}^{-1}$ & $634{,}012~\text{s}^{-1}$ & $25.3~\mu$s & $39{,}526~\text{s}^{-1}$ \\
TON\_IoT        & $148{,}768~\text{s}^{-1}$ & $595{,}072~\text{s}^{-1}$ & $26.9~\mu$s & $37{,}192~\text{s}^{-1}$ \\
MedSec-25       & $134{,}864~\text{s}^{-1}$ & $539{,}456~\text{s}^{-1}$ & $29.7~\mu$s & $33{,}670~\text{s}^{-1}$ \\
WUSTL-EHMS-2020 & $243{,}843~\text{s}^{-1}$ & $975{,}372~\text{s}^{-1}$ & $16.4~\mu$s & $60{,}976~\text{s}^{-1}$ \\
\bottomrule
\end{tabular}
\end{table*}

\subsection{Complete model comparison}
\label{sec:res-complete}

% In your preamble, make sure you have:
% \usepackage{booktabs}
% \usepackage{multirow}
% \usepackage{longtable}

\onecolumn
\begin{longtable}{llcccccr}
\caption{Complete comparison of all models across NSL-KDD, TON-IoT, MedSec-25, and WUSTL-EHMS-2020 datasets.}
\label{tab:complete_results} \\
\toprule
Dataset & Model & Acc & F1 & Prec & Rec & Size (KB) \\
\midrule
\endfirsthead

\multicolumn{7}{c}%
{\tablename\ \thetable\ -- \textit{Continued from previous page}} \\
\toprule
Dataset & Model & Acc & F1 & Prec & Rec & Size (KB) \\
\midrule
\endhead

\midrule
\multicolumn{7}{r}{\textit{Continued on next page}} \\
\endfoot

\bottomrule
\endlastfoot

\multirow{16}{*}{NSL-KDD}
 & XGBoost                & 0.9956 & 0.9488 & 0.9442 & 0.9558 & 1266.1   \\
 & XGBoost\_Cmatched      & 0.9952 & 0.9429 & 0.9419 & 0.9458 & 1161.3   \\
 & GLADE\_FPTM\_final     & 0.9951 & 0.9301 & 0.9399 & 0.9210 & 12.6     \\
 & MLP\_small             & 0.9901 & 0.9159 & 0.9411 & 0.8966 & 115.4    \\
 & MLP\_2C                & 0.9924 & 0.9109 & 0.8983 & 0.9248 & 1205.0   \\
 & MLP\_C                 & 0.9923 & 0.9036 & 0.9115 & 0.8967 & 360.9    \\
 & MLP\_med               & 0.9905 & 0.8986 & 0.9092 & 0.8893 & 242.8    \\
 & MLP\_tiny              & 0.9888 & 0.8907 & 0.9067 & 0.8780 & 58.2     \\
 & kNN\_5                 & 0.9894 & 0.8855 & 0.9052 & 0.8691 & 110463.0 \\
 & LogisticRegression     & 0.9769 & 0.8580 & 0.8877 & 0.8337 & 8.6      \\
 & RandomForest           & 0.9926 & 0.8563 & 0.9647 & 0.8167 & 8490.2   \\
 & LinearSVM              & 0.9655 & 0.8530 & 0.8986 & 0.8194 & 8.6      \\
 & DecisionTree           & 0.9929 & 0.8107 & 0.9261 & 0.7876 & 44.8     \\
 & DecisionTree\_Lmatched & 0.9746 & 0.7277 & 0.7633 & 0.7034 & 11.4     \\
 & RandomForest\_Cmatched & 0.9625 & 0.5843 & 0.5821 & 0.5869 & 961.7    \\
 & NaiveBayes             & 0.7567 & 0.4423 & 0.5765 & 0.6726 & 13.2     \\
\midrule

\multirow{16}{*}{TON\_IoT}
 & XGBoost                & 0.9775 & 0.9511 & 0.9439 & 0.9614 & 2508.2   \\
 & XGBoost\_Cmatched      & 0.9764 & 0.9498 & 0.9424 & 0.9608 & 1899.0   \\
 & RandomForest           & 0.9765 & 0.9467 & 0.9366 & 0.9637 & 14735.4  \\
 & kNN\_5                 & 0.9742 & 0.9441 & 0.9430 & 0.9461 & 51444.2  \\
 & GLADE\_FPTM\_final     & 0.9705 & 0.9418 & 0.9334 & 0.9558 & 16.4     \\
 & DecisionTree           & 0.9508 & 0.9211 & 0.9200 & 0.9301 & 85.5     \\
 & MLP\_2C                & 0.9352 & 0.8830 & 0.9088 & 0.8769 & 914.5    \\
 & MLP\_C                 & 0.9097 & 0.8570 & 0.8802 & 0.8473 & 166.2    \\
 & RandomForest\_Cmatched & 0.8771 & 0.7814 & 0.7961 & 0.7798 & 840.6    \\
 & MLP\_med               & 0.8271 & 0.7734 & 0.8214 & 0.7746 & 125.4    \\
 & MLP\_small             & 0.7983 & 0.7318 & 0.7955 & 0.7238 & 55.6     \\
 & DecisionTree\_Lmatched & 0.7945 & 0.7191 & 0.7985 & 0.7209 & 9.9      \\
 & MLP\_tiny              & 0.7744 & 0.7088 & 0.7482 & 0.7130 & 28.8     \\
 & LogisticRegression     & 0.7076 & 0.6347 & 0.7154 & 0.6494 & 5.5      \\
 & LinearSVM              & 0.7114 & 0.6273 & 0.7106 & 0.6486 & 5.5      \\
 & NaiveBayes             & 0.3875 & 0.3269 & 0.3694 & 0.4201 & 8.5      \\
\midrule

\multirow{16}{*}{MedSec-25}
 & XGBoost                & 0.9768 & 0.8984 & 0.9005 & 0.9004 & 1700.7   \\
 & RandomForest           & 0.9767 & 0.8980 & 0.9022 & 0.9036 & 15734.5  \\
 & XGBoost\_Cmatched      & 0.9753 & 0.8926 & 0.8951 & 0.8951 & 939.0    \\
 & DecisionTree           & 0.9731 & 0.8824 & 0.8953 & 0.8732 & 70.6     \\
 & kNN\_5                 & 0.9733 & 0.8799 & 0.8752 & 0.8848 & 190623.7 \\
 & GLADE\_FPTM\_final     & 0.9707 & 0.8727 & 0.8787 & 0.8828 & 10.5     \\
 & MLP\_2C                & 0.9660 & 0.8657 & 0.8582 & 0.9153 & 702.1    \\
 & MLP\_med               & 0.9680 & 0.8590 & 0.8678 & 0.8621 & 145.3    \\
 & MLP\_C                 & 0.9661 & 0.8453 & 0.8755 & 0.8254 & 161.7    \\
 & MLP\_tiny              & 0.9654 & 0.8440 & 0.8670 & 0.8291 & 33.0     \\
 & RandomForest\_Cmatched & 0.9611 & 0.8390 & 0.8631 & 0.8367 & 524.7    \\
 & MLP\_small             & 0.9646 & 0.8372 & 0.8630 & 0.8198 & 65.8     \\
 & LogisticRegression     & 0.9413 & 0.7763 & 0.7891 & 0.7853 & 4.8      \\
 & DecisionTree\_Lmatched & 0.9392 & 0.7203 & 0.8067 & 0.6782 & 7.3      \\
 & LinearSVM              & 0.8538 & 0.5146 & 0.6990 & 0.4660 & 4.7      \\
 & NaiveBayes             & 0.7673 & 0.4218 & 0.4803 & 0.5192 & 6.8      \\
\midrule

\multirow{16}{*}{WUSTL-EHMS-2020}
 & MLP\_2C                & 0.9485 & 0.8248 & 0.8997 & 0.7917 & 349.8   \\
 & GLADE\_FPTM\_final     & 0.9442 & 0.7976 & 0.9050 & 0.7645 & 6.2     \\
 & kNN\_5                 & 0.9375 & 0.7922 & 0.8445 & 0.7684 & 2857.4  \\
 & XGBoost                & 0.9406 & 0.7590 & 0.9151 & 0.7314 & 615.3   \\
 & MLP\_C                 & 0.9375 & 0.7465 & 0.8774 & 0.7237 & 70.3    \\
 & DecisionTree           & 0.9357 & 0.7216 & 0.8902 & 0.7049 & 19.7    \\
 & XGBoost\_Cmatched      & 0.9354 & 0.7118 & 0.9198 & 0.6994 & 388.0   \\
 & DecisionTree\_Lmatched & 0.9295 & 0.6936 & 0.7943 & 0.6893 & 7.1     \\
 & RandomForest           & 0.9320 & 0.6778 & 0.8910 & 0.6807 & 3562.2  \\
 & MLP\_med               & 0.9298 & 0.6603 & 0.9066 & 0.6713 & 102.1   \\
 & RandomForest\_Cmatched & 0.9305 & 0.6554 & 0.9671 & 0.6692 & 380.6   \\
 & MLP\_tiny              & 0.9289 & 0.6507 & 0.9602 & 0.6669 & 20.7    \\
 & MLP\_small             & 0.9280 & 0.6497 & 0.7572 & 0.6682 & 43.5    \\
 & LinearSVM              & 0.9286 & 0.6449 & 0.6253 & 0.6657 & 2.4     \\
 & LogisticRegression     & 0.9280 & 0.6430 & 0.6220 & 0.6655 & 2.5     \\
 & NaiveBayes             & 0.1293 & 0.3643 & 0.6570 & 0.6668 & 3.0     \\

\end{longtable}
\twocolumn
%% ========================================================================
\begin{thebibliography}{99}

\bibitem{granmo2018tsetlin}
O.-C. Granmo,
``The Tsetlin machine --- a game theoretic bandit driven approach to
optimal pattern recognition with propositional logic,''
\emph{arXiv preprint arXiv:1804.01508}, 2018.

\bibitem{hnilov2025fuzzy}
A. Hnilov,
``Fuzzy-pattern Tsetlin machine,''
\emph{arXiv preprint arXiv:2508.08350}, 2025.

\bibitem{granmo2019convolutional}
O.-C. Granmo, S. Berge, L. Jiao, X. Zhang, G. T. Berge and M. Goodwin,
``The convolutional Tsetlin machine,''
\emph{arXiv preprint arXiv:1905.09688}, 2019.

\bibitem{abeyrathna2019continuous}
K. D. Abeyrathna, O.-C. Granmo, X. Zhang and M. Goodwin,
``A scheme for continuous input to the Tsetlin machine with
applications to forecasting disease outbreaks,''
\emph{arXiv preprint arXiv:1905.04199}, 2019.

\bibitem{abeyrathna2021adaptive}
K. D. Abeyrathna, O.-C. Granmo, X. Zhang, L. Jiao and M. Goodwin,
``Adaptive sparse representation of continuous input for Tsetlin
machines based on stochastic searching on the line,''
\emph{Electronics}, vol. 10, no. 17, art. 2107, 2021.

\bibitem{maheshwari2023redress}
S. Maheshwari, T. Rahman, R. Shafik, A. Yakovlev, A. Rafiev, L. Jiao and
O.-C. Granmo,
``REDRESS: Generating compressed models for edge inference using
Tsetlin machines,''
\emph{IEEE Trans. Pattern Anal. Mach. Intell.}, vol. 45, no. 9,
pp. 11077--11091, Sep. 2023.

\bibitem{bhattarai2023concise}
B. Bhattarai, O.-C. Granmo and L. Jiao,
``Building concise logical patterns by constraining Tsetlin machine
clause size,''
in \emph{Proc. 32nd Int. Joint Conf. Artif. Intell. (IJCAI)},
Macao, China, Aug. 2023, pp. 3391--3399.

\bibitem{sharma2023drop}
J. Sharma, R. Yadav, O.-C. Granmo and L. Jiao,
``Drop clause: Enhancing performance, interpretability and robustness
of the Tsetlin machine,''
in \emph{Proc. 37th AAAI Conf. Artif. Intell.}, Washington, DC, USA,
Feb. 2023, pp. 13547--13555.

\bibitem{tavallaee2009nslkdd}
M. Tavallaee, E. Bagheri, W. Lu and A. A. Ghorbani,
``A detailed analysis of the KDD CUP 99 data set,''
in \emph{Proc. IEEE Symp. Comput. Intell. Secur. Defense Appl. (CISDA)},
Ottawa, ON, Canada, Jul. 2009, pp. 1--6.

\bibitem{moustafa2021ton}
N. Moustafa,
``A new distributed architecture for evaluating AI-based security
systems at the edge: Network TON\_IoT datasets,''
\emph{Sustainable Cities and Society}, vol. 72, art. 102994, Sep.
2021.

\bibitem{hady2020wustl}
A. A. Hady, A. Ghubaish, T. Salman, D. Unal and R. Jain,
``Intrusion detection system for healthcare systems using medical and
network data: A comparison study,''
\emph{IEEE Access}, vol. 8, pp. 106576--106584, 2020.

\bibitem{chen2016xgboost}
T. Chen and C. Guestrin,
``XGBoost: A scalable tree boosting system,''
in \emph{Proc. 22nd ACM SIGKDD Int. Conf. Knowl. Discov. Data Min.},
San Francisco, CA, USA, Aug. 2016, pp. 785--794.

\bibitem{breiman2001random}
L. Breiman,
``Random forests,''
\emph{Machine Learning}, vol. 45, no. 1, pp. 5--32, Oct. 2001.

\bibitem{pedregosa2011scikit}
F. Pedregosa \emph{et al.},
``Scikit-learn: Machine learning in Python,''
\emph{Journal of Machine Learning Research}, vol. 12, pp. 2825--2830,
2011.

\end{thebibliography}

\end{document}
